\section{Deployment Recommendation and Remaining Limits}
\paragraph{Choose a service before an algorithm.}
Do not deploy the strict wrapper solely on the strength of these transmission savings. For local trusted status, use the root projection; for final per-test states, a flat ledger may suffice. Where protected hierarchical evidence is already required, enable a suitable codec, retain the incumbent, and consider learned-order balanced fixed8 as the low-fitting-cost candidate. Exact DP remains an offline reference and can be worthwhile when its extra byte reduction pays back. Block DP is an ablation, not the default.

\paragraph{A cost calculator is not measured payback.}
The new scenarios expose when assumed extra costs dominate compressed delivery savings. They neither import a five-second overhead from a different fixture as though paired with each source case nor assume unlimited repeat exports. Dynamic adaptation is optional; the priced window selects already compiled candidates and can regress under rapid change. Whole-epoch compression changes streaming latency, and high-latency links do not necessarily favor HACT. Those are deployment questions, not consequences of the tree objective.

\paragraph{New default, new evidence boundary.}
The practical recommendation changed after the v4 review. It is therefore tested on new frozen synthetic seeds with disjoint train/validation/held-out draws. The twelve source candidates are reused diagnostic evidence and are not an independent test of model-generated repairs. On these replays, learned balanced trees are smaller than original-order fixed8 yet larger than exact learned layouts. The recommendation is about compilation simplicity, not universal superiority. Three seeds per setting support descriptive ranges, not precise generalization claims.

\paragraph{Unclosed empirical requests.}
Neither a physical WAN nor actual public-vocabulary BPE was evaluated in this release. Software/vocabulary retrieval failed, and no authorized pair of remote controlled hosts was available. The optional tools and local tests do not close these experimental gaps. No hosted model was called, and there is no new SWE-bench or natural historical-bug cohort. The actual SWE-bench harness uses isolated repository environments~\cite{sweharness}; scripted source mutations are not equivalent. Source-level token counts, when measured later, will still not establish provider billing or better agent reasoning.

\paragraph{Trust and scope.}
Future Code remains a single-host supervisor. The guard and TLS utility are not a distributed consensus protocol or fault-tolerant deployment. A passing declared suite cannot establish absence of arbitrary defects. Authorization requires an uncompromised checker, correctly scoped inputs and protected reference state. A separate integration service must bind the verified candidate and integration head atomically; no PDF table or root JSON card substitutes for that step.

\section{Conclusion}
Adaptive HACT preserves an acceptance contract while changing the organization of its evidence. Full completion hyperedges supply the cost signal; stable IDs and generation fencing preserve local meaning. The practical path now uses a balanced tree over a learned order, retains the incumbent and performs no DP. New held-out synthetic tests and archived-source replays support a limited byte-efficiency claim, not universal latency gains. Service-matched accounting makes the adoption boundary explicit, while the missing physical-WAN, BPE and live-agent experiments remain identified rather than replaced by proxies. The result is a simpler implementation and a more falsifiable deployment claim.
